\section{Higher quotients: keeping the identifications}
\label{h:chapter}

The exact Mal'cev argument has answered a specific question: which simple
quotient objects survive every choice of finite normal series?  We now ask
what was omitted when we retained only those quotient objects.  A quotient
can remember not only which points have been identified, but also the ways
of identifying them and the compatibility between successive identifications.
This is where the higher framework becomes useful.

The change of foundations in this section is explicit.  We use spaces,
homotopy limits, and a small number of stated theorems about
$\infty$-categories.  Their foundations are not proved here.  A result marked
\emph{Framework theorem, quoted} is an external input; the propositions
following it are deductions with their own proofs.  None of these inputs
is used in the earlier proof of Jordan--H\"older.  Conversely, invoking the
word ``higher'' will not be used to justify an unproved quotient law.

\subsection{One construction of homotopical algebraic objects}

Write $\mathcal S$ for the $\infty$-category of spaces, or homotopy types.
A set is a discrete space.  A space can also have paths, homotopies between
paths, and further coherences.  The mapping object
$\operatorname{Map}(X,Y)$ is consequently a space: its points are maps,
its paths are homotopies, and so on.  An equivalence is the appropriate
replacement for an isomorphism.  Limits and colimits in this setting mean
homotopy-coherent limits and colimits.

An ordinary one-sorted algebraic theory $\mathbb T$ is a small category
with finite products, generated under products by an object $t$.
The objects $t^n$ stand for $n$ inputs.  An arrow $t^n\to t$ stands for
an $n$-ary operation, and composition records substitution.  Equations
between arrows record equations between operations.  A model in a category
with finite products is a finite-product-preserving functor out of
$\mathbb T$.

For groups, multiplication is an arrow $t^2\to t$, inverse is $t\to t$,
and identity is $t^0\to t$.  Associativity is equality of the two arrows
$t^3\to t$ obtained by substitution.  For left $R$-modules we include
addition, zero, inverse, and one unary operation for each $r\in R$.
The equation for multiplication of scalars is an equation between
composites of these unary operations.  Thus the presentation does not
discard the ring when it passes to a common language.

Define
\begin{equation}\label{h:theory-models}
 \mathcal C_{\mathbb T}
   =\operatorname{Fun}^{\times}(\mathbb T,\mathcal S).
\end{equation}
Preservation of products now means preservation up to the specified coherent
equivalences.  This is the promised single construction: change the target
of the theory from sets to spaces, retaining the operations and their
compatibilities.

\begin{proposition}[Recovering the ordinary algebra]\label{h:discrete}
The full subcategory of $0$-truncated objects of $\mathcal C_{\mathbb T}$
is equivalent to the nerve of the ordinary category
$\operatorname{Fun}^{\times}(\mathbb T,\mathbf{Set})$.
\end{proposition}
\begin{proof}
An object $M$ is $0$-truncated when $\operatorname{Map}(P,M)$ is discrete
for every $P$.  If all $M(t^n)$ are discrete, the space of natural
transformations into $M$ is a limit of discrete mapping spaces and is
therefore discrete.  Thus such $M$ is $0$-truncated.

Conversely, the covariant representable
$F=\operatorname{Hom}_{\mathbb T}(t,-)$ preserves products, so belongs to
$\mathcal C_{\mathbb T}$.  Yoneda identifies
$\operatorname{Map}(F,M)$ with $M(t)$: a natural transformation is
determined by its value on $\id_t$, and the value on an arrow
$a:t\to t^n$ is then $M(a)$ applied to that value. This is the
space-valued Yoneda lemma, including its coherent mapping-space statement.
If $M$ is $0$-truncated, $M(t)$
is discrete; product preservation gives $M(t^n)\simeq M(t)^n$, so all its
values are discrete.  Functors between ordinary discrete values have
ordinary natural transformations and the ordinary product-preservation
condition.  This identifies both objects and maps of the displayed
subcategory.
\end{proof}

\begin{theorem}[Framework theorem, quoted: rigidification]
\label{h:rigidification}
For a small ordinary finite-product theory $\mathbb T$, the homotopy theory
of strict simplicial $\mathbb T$-models, with objectwise weak equivalences,
presents $\mathcal C_{\mathbb T}$.
\end{theorem}
This is Lurie's Corollary 5.5.9.3, using Proposition 5.5.9.2, in
\cite{HTT}; see also \cite{Badzioch}.  This theorem is what permits us
to compute with familiar simplicial algebra while interpreting the answer
in \eqref{h:theory-models}.  It is not a claim that all coherences can
simply be omitted.

For group theory, the result is the homotopy theory of simplicial groups.
Two further foundational inputs describe it. The Kan loop-group and
classifying-space adjunction identifies this homotopy theory with pointed
connected spaces \cite[Chapter V, Proposition 6.3 and Corollary 6.4]{GoerssJardine}.
The loop--delooping recognition theorem identifies pointed connected spaces
with grouplike $E_1$-spaces \cite[Theorem 5.2.6.10]{HA}.
An $E_1$ multiplication is associative with its coherent homotopies;
``grouplike'' says that its components form a group. The two quoted
equivalences together supply the asserted identification; rigidification
alone does not supply either comparison.

For module theory, the Dold--Kan theorem identifies simplicial $R$-modules
with nonnegative chain complexes of $R$-modules; weak equivalences become
quasi-isomorphisms \cite[Tag 019G]{Stacks}.  Localizing therefore gives
connective derived $R$-modules.  For a commutative ring $R$, this agrees
with connective modules over its Eilenberg--Mac Lane ring spectrum $HR$
\cite[Theorem 7.1.2.13]{HA}.  The $HR$-module structure is essential:
an arbitrary spectrum, even a connective one, is not thereby an
$R$-module.

The earlier theorem applies to the ordinary models recovered by
Proposition~\ref{h:discrete}, whenever their category satisfies its pointed
exact Mal'cev hypotheses.  Groups and modules are therefore instances
inside a common homotopical construction.  The exactness hypotheses do
real work; product-preserving models in spaces alone do not imply a
composition theorem for every higher object.

\subsection{The action groupoid and the role of associativity}

Let an ordinary group $G$ act on the right on a set $X$.  Form a groupoid
whose objects are elements $x\in X$, and whose arrows are
\[
 (x,g):x\longrightarrow xg.
\]
Composition is
\[
 (x,g)\text{ followed by }(xg,h)=(x,gh).
\]
It is defined because $(xg)h=x(gh)$; it is associative because multiplication
in $G$ is associative; identities and inverses follow from the corresponding
group laws.  This construction retains the actual witness $g$ of an
identification, not only the fact that two points lie in the same orbit.

Its nerve is the simplicial set
\begin{equation}\label{h:action-nerve}
 X_n=X\times G^n.
\end{equation}
In degree $n$, the tuple $(x,g_1,\ldots,g_n)$ means the composable chain
\[
 x\longrightarrow xg_1\longrightarrow xg_1g_2
   \longrightarrow\cdots\longrightarrow xg_1\cdots g_n.
\]
The first face removes $x$ by replacing it with $xg_1$; an interior face
replaces $g_i,g_{i+1}$ by their product; the last face removes $g_n$.
A degeneracy inserts an identity element.  For three consecutive arrows,
the equality of the two ways to take interior faces is precisely
$(g_1g_2)g_3=g_1(g_2g_3)$.  The unit identities similarly supply the
face--degeneracy identities.  The nerve is thus a way to organize, in
one object, the algebraic compatibility behind the elementary maneuvers.

\begin{definition}\label{h:homotopy-quotient}
The homotopy quotient $X/\!/G$ is the space obtained by realizing the
nerve \eqref{h:action-nerve}.  For a point with the trivial action, write
$BG=*/\!/G$.
\end{definition}

\begin{proposition}[What an action quotient retains]\label{h:stabilizers}
The connected components of $X/\!/G$ are the orbits of $X$.  The component
of $x$ is equivalent to $BG_x$, where
$G_x=\{g\in G:xg=x\}$.  In particular a free action has homotopy quotient
equivalent to its discrete orbit set.
\end{proposition}
\begin{proof}
There is an arrow from $x$ to $y$ exactly when $y=xg$ for some $g$.
Thus components are orbits.  In a fixed orbit, the full subgroupoid on
one representative $x$ has automorphism group $G_x$.  Its inclusion is
fully faithful by construction and essentially surjective because every
object in the orbit is isomorphic to $x$.  It is therefore an equivalence
of groupoids.  Realization preserves such equivalences, giving $BG_x$.
For a free action all stabilizers are trivial, so each component is
contractible.
\end{proof}

For the trivial action on a point, the ordinary orbit set is a point,
whereas the homotopy quotient is $BG$ and retains every automorphism in
$G$.  For a subgroup $N\leq G$ acting on $G$ by right translation, the
action is free, so $G/\!/N\simeq G/N$ as spaces.  If $N$ is normal,
multiplication also descends to these cosets.  The two assertions are
different: freeness determines the homotopy type of the quotient;
normality determines its induced group multiplication.

\subsection{Extensions become fiber diagrams}

An ordinary pullback asks for equal images. A homotopy pullback
$X\times^h_ZY$ retains a path between the images: a point is a pair
$(x,y)$ together with a path between $f(x)$ and $g(y)$.
Paths between such data must respect the specified comparison, with the
corresponding higher coherences. In particular, a homotopy fiber over
$z$ remembers both a point $x$ and a path between $f(x)$ and $z$.
The next calculation is a case where this extra datum is a group element.

\begin{proposition}\label{h:extension-fiber}
A short exact sequence of ordinary groups
\[
 1\longrightarrow N\longrightarrow G\xrightarrow{p}Q\longrightarrow1
\]
determines a homotopy pullback square
\[
\begin{tikzcd}
 BN \arrow[r] \arrow[d] & BG \arrow[d,"Bp"]\\
 * \arrow[r] & BQ.
\end{tikzcd}
\]
\end{proposition}
\begin{proof}
The homotopy fiber of the functor between one-object groupoids induced by
$p$ can be modeled by the groupoid with objects $q\in Q$ and arrows
$g:q\to q p(g)$.  Indeed its objects are an object of $BG$ together with
an isomorphism from the chosen object of $BQ$ to its image, and its arrows
are the compatible isomorphisms.  This is the groupoid model of the
homotopy pullback, retaining the isomorphism in the pullback datum.
Surjectivity of $p$ makes it transitive, and the automorphisms of $1\in Q$
are exactly $\ker p=N$.  The argument of
Proposition~\ref{h:stabilizers} identifies this fiber with $BN$.
The inclusion of its representative object gives the usual map $BN\to BG$.
\end{proof}

This square is a fiber description, not a claim about an ordinary
set-theoretic fiber of the map of realizations.  It also is not automatically
a cofiber square.  The distinctions matter because the homotopy theory
of spaces is unstable.

\subsection{Quotients that survive change of base}

An $\infty$-topos is a setting for spaces varying over a site; formally it
is an accessible left exact localization of a presheaf $\infty$-category
of spaces.  The theory-model category $\mathcal C_{\mathbb T}$ need not
be an $\infty$-topos.  The effective-quotient results below apply in spaces
or another specified $\infty$-topos; the subsequent descent equivalence
is stated for spaces.  We need the following two consequences of the
foundational theory of an $\infty$-topos.

\begin{theorem}[Framework theorem, quoted: effective quotients]
\label{h:topos-foundation}
In an $\infty$-topos, colimits are preserved by pullback.  Every groupoid
object is the \v{C}ech nerve of the map from its object of objects to its
realization.
\end{theorem}
These are the universality and effectivity assertions in
\cite[Theorem 6.1.0.6]{HTT}.  A \v{C}ech nerve of $U\to Y$ has degree
$n$ term $U\times_Y\cdots\times_Y U$ with $n+1$ factors.  Effectivity
says that the identification data in a groupoid object are precisely the
data recovered by pulling back its quotient map against itself.

\begin{proposition}[A quotient can be computed after base change]
\label{h:quotient-basechange}
In an $\infty$-topos, let $U_\bullet$ be a groupoid object with realization $Y$, and let
$f:Y'\to Y$.  Then
\[
 \left|Y'\times_Y U_\bullet\right|\simeq Y'.
\]
Moreover $Y'\times_Y U_\bullet$ is the \v{C}ech nerve of
$Y'\times_Y U_0\to Y'$.
\end{proposition}
\begin{proof}
Pullback commutes with the realization colimit by
Theorem~\ref{h:topos-foundation}, so its realization is
$Y'\times_Y|U_\bullet|\simeq Y'\times_Y Y\simeq Y'$.
Effectivity identifies $U_n$ with the $(n+1)$-fold fiber product of $U_0$
over $Y$.  Pulling back this description gives the corresponding
$(n+1)$-fold fiber product of $Y'\times_Y U_0$ over $Y'$, including its
face and degeneracy maps.  These identifications prove the nerve assertion.
\end{proof}

The usefulness is concrete.  A quotient description can be transported to
a new parameter space without recomputing its coherence relations.  If we
make two successive base changes, the two answers agree by the canonical
associativity equivalence of fiber products.  That equivalence is itself
compatible with a third base change: both descriptions are limits of the
same diagram.  Universal properties account for the compatibility that
would otherwise have to be checked separately in each example.

\subsection{Descent: a family over \texorpdfstring{$BG$}{BG} is an action}

\begin{theorem}[Framework theorem, quoted: straightening for spaces]
\label{h:straightening}
For a space $B$, taking homotopy fibers with transport gives an equivalence
\[
 \mathcal S_{/B}\simeq\operatorname{Fun}(B,\mathcal S).
\]
\end{theorem}
This is the groupoid case of straightening and unstraightening
\cite[\S\S2.2,3.2]{HTT}.  The symbol $B$ on the right is regarded as
an $\infty$-groupoid.  Transport along an invertible path is an equivalence;
transport along composable paths includes coherent composition data.

For $B=BG$, a functor $BG\to\mathcal S$ is a space with coherent
$G$-action.  A map $E\to BG$ determines the fiber $X$ over the basepoint,
and loops in $BG$ transport $X$ around the base.  Conversely, the
Grothendieck construction assembles the action into its total space
$X/\!/G\to BG$.  These are the actual forward and inverse constructions
in this instance of the theorem.  In identifying this functor notation
with the usual left-action convention, we use the canonically equivalent
opposite one-object model via inversion: the earlier right action becomes
$g\cdot x=xg^{-1}$.  Our concrete groupoid calculations continue to use
the right-multiplication convention, with first $g$, then $h$, composed
as $gh$.  No change of coefficient structure is involved.

\begin{corollary}[Equivariant information is descent information]
\label{h:action-descent}
The equivalence
\[
 \mathcal S_{/BG}\simeq\operatorname{Fun}(BG,\mathcal S)
\]
identifies maps over $BG$ with coherent equivariant maps.  Its restriction
along $BH\to BG$ for a subgroup $H$ identifies homotopy pullback of families
with restriction of their actions.  Sections of $X/\!/G\to BG$ form the
space of coherent fixed points of $X$.
\end{corollary}
\begin{proof}
The first assertion is the assertion about mapping spaces in the
equivalence.  Taking the fiber of a pulled-back family at a point of $BH$
gives the same fiber as before, and its path transport is the original
transport along the image path in $BG$.  This is precisely restriction
along the functor $BH\to BG$.  Finally a section is a map over $BG$ from
the identity family $BG\to BG$.  That family corresponds to the constant
one-point functor.  Maps from it to the action functor are compatible
choices of a point, with the homotopies and coherences making it fixed.
\end{proof}

For a discrete $G$-set these are ordinary fixed points.  For a general
space the coherences are part of the answer.  We have gained a theory of
families, transport, and descent, not merely another notation for a set
of orbits.

\subsection{Intersections acquire additional components}

Here is a useful calculation that prevents a misleading translation of
the earlier subgroup diagrams.

\begin{proposition}[Double cosets in a homotopy pullback]
\label{h:double-coset}
For subgroups $A,B\leq G$, there is an equivalence of spaces
\[
 BA\times^h_{BG}BB
 \simeq
 \coprod_{[g]\in A\backslash G/B}
 B(A\cap gBg^{-1}).
\]
The component selected by the identity element is $B(A\cap B)$.
\end{proposition}
\begin{proof}
A groupoid model for the homotopy pullback has one object for each
$g\in G$.  An arrow from $g$ to $g'$ is a pair $(a,b)\in A\times B$
such that $g'=a^{-1}gb$.  First $(a,b)$ and then $(a',b')$ compose to
$(aa',bb')$, since $a'^{-1}(a^{-1}gb)b'=(aa')^{-1}g(bb')$.
This is the weak pullback: $g$ records the comparison isomorphism in
$BG$ between the images of the two chosen objects.

Its orbits are the double cosets $AgB$.  The stabilizer of $g$ consists
of pairs with $a=gbg^{-1}$ and is identified with $A\cap gBg^{-1}$ by
projection to $a$.  Apply the transitive-groupoid argument of
Proposition~\ref{h:stabilizers} on each orbit.  For $g=1$ the stabilizer
is $A\cap B$, as claimed.
\end{proof}

The classical intersection is present, but it is one component of a richer
comparison object.  Discarding the other components before doing base
change would discard genuine information.  This is one reason to keep
pointings and transport explicit when moving a subgroup argument into
homotopy theory.

In particular, a proper subgroup inclusion $H\hookrightarrow G$ does not
give a monomorphism $BH\to BG$ in $\mathcal S$.  Its homotopy fiber is the
coset space, which has more than one point.  A monomorphism of spaces has
every fiber either empty or contractible, so this map is not one.
The subgroup inclusion is monic in the category of discrete groups;
its delooping has a different categorical type.

\subsection{Higher normality requires quotient data}

For a higher group map $H\to G$, a normality datum can be specified by a
pointed connected space $Q$ and a homotopy fiber square in the ambient
$\infty$-category of spaces
\[
\begin{tikzcd}
 BH \arrow[r] \arrow[d] & BG \arrow[d]\\
 * \arrow[r] & Q.
\end{tikzcd}
\]
All displayed maps preserve the specified points. Requiring a pullback
in spaces retains the entire homotopy fiber; taking a limit only inside
the subcategory of pointed connected spaces could discard other components.
The loop space $\Omega Q$ then plays the role of the quotient higher
group.  For an ordinary normal inclusion, $Q=B(G/H)$ supplies this datum.
For general higher maps, the datum is additional structure and need not
be unique.

The appropriate isomorphism laws consequently include this structure.
Beardsley and Fox develop normality data for group objects in
$\infty$-topoi.  Their current formulation proves the first and second
Noether theorems, exhibits failures of an unrestricted third theorem,
and proves an iterated-quotient theorem with compatible data
\cite[\S5, especially Theorem 5.14]{BeardsleyFox}.  We import no general
higher Jordan--H\"older theorem from that work.

One can already see why the ordinary definition of a filtration cannot
be transferred without a decision about admissible arrows.

\begin{proposition}[Literal subobject chains force discreteness]
\label{h:monos-discrete}
Suppose a group object in spaces admits a finite chain of categorical
monomorphisms beginning at the trivial group and ending at $G$.  Then
the underlying space of $G$ is discrete up to equivalence.
\end{proposition}
\begin{proof}
Limits of product-preserving theory models are computed pointwise,
because limits commute with products; evaluation on the generator
detects equivalences because its powers give every object of the theory.
Thus limits of group objects are computed on their underlying spaces,
and equivalences are detected there.  A map is monic precisely when its
diagonal is an equivalence.  Hence the forgetful functor detects these
monomorphisms.  Their composite gives a monomorphism $*\to G$ of spaces.
Its fiber over the identity is $\Omega G$, so that loop space is
contractible.  The identity component of $G$ is therefore contractible.
Translation by a point of a group object identifies its component with
the identity component.  All components are contractible, which is
exactly the assertion that $G$ is $0$-truncated.
\end{proof}

Thus a theory of genuinely higher composition objects must specify more
than the old finite chain of monomorphisms.  Our common ordinary theorem
is complete in its stated exact framework.  The higher enlargement gives
coherent quotients, families, and descent; any additional composition
theorem must state and verify the admissible quotient and refinement
conditions it uses.  This is a precise boundary of the proved result,
not a reason to discard the enlargement.
